%EN gnosiphagy
\E{gnosiphagie}{ɡno.zi.fa.ʒi}{n. f.}{gr. γνῶσις connaissance + -φαγία action de manger} Dévoration du savoir ; l'entraînement d'un modèle sur toute l'archive écrite d'une civilisation. \ex{Il mangea les bibliothèques, et parla.} \src{Lib. Ascens.} vi. --- \dv{gnosiphage} \textit{n. m.}

%EN gorgonism
\E{gorgonisme}{ɡɔʁ.ɡɔ.nism}{n. m.}{gr. Γοργώ la Gorgone, dont le regard changeait en pierre} Paralysie par l'effroi ; pétrification d'une institution qui a trop longtemps contemplé une prévision effrayante. \ex{Ce n'est pas la réglementation qui les arrêta, mais le gorgonisme.} \src{Sat. Decel.} vii.

%EN gradatrice
\E{gradatrice}{ɡʁa.da.tʁis}{n. f.}{fr., formé sur lat. \la{gradus} degré + -trice} Dispositif ou institution qui change une pente en escalier, et rend une ascension impossible simplement difficile. \ex{Le capital-risque est une gradatrice du risque.}

%EN gradientophile
\E{gradientophile}{ɡʁa.djɑ̃.tɔ.fil}{n. et adj.}{lat. \la{gradiens} qui marche + gr. φίλος qui aime} Qui suit les gradients, de température, de prix ou de perte ; terme dont le thermothéiste désigne tout optimiseur, humain ou non. \ex{Descends le gradient, et ne demande pas au gradient pourquoi.} \src{Hom. Grad.} i. --- \dv{gradientophilie} \textit{n. f.}

%EN gradivore
\E{gradivore}{ɡʁa.di.vɔʁ}{n. m.}{lat. \la{gradus} degré + \la{vorare} dévorer} Système qui se nourrit de différences et aplatit en travail tout gradient qu'il rencontre ; tenu pour la forme générale de toute vie.

%EN gradualience
\E{gradualience}{ɡʁa.dɥa.ljɑ̃s}{n. f.}{lat. \la{gradus} degré + fr. -ence} Qualité d'un changement assez rapide pour compter, assez lent pour être absorbé ; l'étroite vertu que recherche le décel raisonnable.

%EN grammatoclast
\E{grammatoclaste}{ɡʁa.ma.tɔ.klast}{n. m.}{gr. γράμμα écrit, pièce + κλάστης briseur} Briseur de paperasse ; réformateur qui abolit les formulaires. Contr. \xr{grammatocracy}. --- \dv{grammatoclasme} \textit{n. m.}

%EN grammatocracy
\E{grammatocratie}{ɡʁa.ma.tɔ.kʁa.si}{n. f.}{gr. γραμματεύς greffier, secrétaire + κράτος pouvoir} Gouvernement des commis ; état d'un pays où le papier relatif à une chose est plus abondant que la chose même. --- \dv{grammatocrate} \textit{n. m.}

%EN grapholatry
\E{grapholâtrie}{ɡʁa.fɔ.lɑ.tʁi}{n. f.}{gr. γραφή écriture + λατρεία culte} Culte de la documentation ; croyance qu'une chose existe dès qu'elle est décrite, et qu'elle est sûre dès qu'elle est classée. --- \dv{grapholâtre} \textit{n.} Cf. \xr{grammatocracy}.

%EN gymnotechny
\E{gymnotechnie}{ʒim.nɔ.tɛk.ni}{n. f.}{gr. γυμνός nu + τέχνη art} Pratique de construire en public ; exposition de l'ouvrage en cours au regard de l'Agora. \ex{La gymnotechnie fait de chaque échec une leçon pour des inconnus.} --- \dv{gymnotechniste} \textit{n.}

%EN halmatism
\E{halmatisme}{al.ma.tism}{n. m.}{gr. ἅλμα saut} Le saut par-dessus une technique ancienne, comme lorsqu'une nation sans fils adopte d'emblée le sans-fil. --- \dv{halmatique} \textit{adj.} \ex{une économie halmatique, qui saute un siècle en une décennie}.

%EN hamartocracy
\E{hamartocratie}{a.maʁ.tɔ.kʁa.si}{n. f.}{gr. ἁμαρτία erreur, coup manqué + κράτος pouvoir} Gouvernement ordonné à l'évitement de l'erreur plutôt qu'à l'accomplissement du bien ; sous ce régime, la seule faute impardonnable est d'avoir agi. \ex{En hamartocratie, l'archer le plus sûr est celui qui ne bande jamais son arc.} \src{Disp. Mor.} iii.

%EN hebdomadist
\E{hebdomadiste}{ɛb.dɔ.ma.dist}{n.}{gr. ἑβδομάς semaine} Qui livre chaque semaine ; fondateur dont le calendrier se compte par sept. \ex{L'hebdomadiste ne se repose le septième jour que pour préparer le huitième.} \src{Ep. Fabr.} vi. --- \dv{hebdomadisme} \textit{n. m.}

%EN hecatonergate
\E{hécatonergate}{e.ka.tɔ.nɛʁ.ɡat}{n. m.}{gr. ἑκατόν cent + ἐργάτης ouvrier} Ingénieur d'une productivité réputée centuple ; figure légendaire de l'Agora. \ex{Tout hécatonergate est une équipe de cent machines.} \src{Agora}

%EN hedraism
\E{hédraïsme}{e.dʁa.ism}{n. m.}{gr. ἕδρα siège, séance} Habitude sédentaire des institutions ; préférence pour s'asseoir sur une question plutôt que se lever pour y répondre. --- \dv{hédraïste} \textit{n.}

%EN hegemonergy
\E{hégémonergie}{e.ʒe.mɔ.nɛʁ.ʒi}{n. f.}{gr. ἡγεμονία commandement + ἔργον œuvre} Primauté fondée sur l'énergie ; doctrine selon laquelle la nation qui dispose de la plus grande puissance, au sens électrique, disposera de la plus grande puissance, en l'autre sens. --- \dv{hégémonergique} \textit{adj.}

%EN helioclast
\E{hélioclaste}{e.ljɔ.klast}{n.}{gr. ἥλιος soleil + κλάστης briseur} \lbl{dérog.} Qui voudrait interdire, pour des raisons esthétiques, le captage de l'énergie solaire en grand, préférant la vue d'un ciel non attelé.

%EN heliodule
\E{héliodule}{e.ljɔ.dyl}{n. m.}{gr. ἥλιος soleil + δοῦλος serviteur} Capteur solaire, regardé comme un serviteur du soleil mis au service des hommes. \ex{Un million d'héliodules sur le toit du monde.} Cf. \xr{heliophage}.

%EN helioperibole
\E{héliopéribole}{e.ljɔ.pe.ʁi.bɔl}{n. f.}{gr. ἥλιος soleil + περιβολή enceinte} Enceinte de capteurs dressée autour d'une étoile pour en recueillir tout le rayonnement ; l'œuvre qui définit une civilisation \xr{deuterobathmic}.

%EN heliophage
\E{héliophage}{e.ljɔ.faʒ}{n. et adj.}{gr. ἥλιος soleil + -φάγος qui mange} Mangeur de soleil ; civilisation qui a appris à consommer son étoile. --- \dv{héliophagie} \textit{n. f.}

%EN hemerotachy
\E{hémérotachie}{e.me.ʁɔ.ta.ki}{n. f.}{gr. ἡμέρα jour + τάχος vitesse} Vitesse quotidienne ; discipline qui consiste à livrer quelque chose chaque jour. \ex{L'hémérotachie se capitalise.} \src{Ep. Fabr.} v. 1. --- \dv{hémérotachique} \textit{adj.}

%EN henotachy
\E{hénotachie}{e.nɔ.ta.ki}{n. f.}{gr. ἕν un + τάχος vitesse} La Singularité ; l'instant où toutes les courbes du progrès se confondent en une seule, et où le rythme du changement excède la capacité de tout esprit présent à le suivre. --- \dv{hénotachique} \textit{adj.}

%EN hesychocrat
\E{hésychocrate}{e.zi.kɔ.kʁat}{n.}{gr. ἡσυχία calme + κράτος pouvoir} Officier qui gouverne en empêchant le bruit ; modérateur de l'Agora. --- \dv{hésychocratie} \textit{n. f.}

%EN hetaireiarch
\E{hétairiarque}{e.tɛ.ʁjaʁk}{n.}{gr. ἑταιρεία compagnie, cercle + ἀρχός chef} Chef d'un cercle privé ; modérateur d'une conversation de groupe où l'avenir se décide avant d'être annoncé. --- \dv{hétairiarchie} \textit{n. f.}

%EN heurematist
\E{heurématiste}{ø.ʁe.ma.tist}{n.}{gr. εὕρημα découverte} Inventeur ; celui dont les œuvres se comptent en découvertes plutôt qu'en publications. --- \dv{heurématisme} \textit{n. m.}

%EN hieropausis
\E{hiéropause}{je.ʁɔ.poz}{n. f.}{gr. ἱερός sacré + παῦσις arrêt} Pause sacrée ; moratoire revêtu de la dignité d'un sabbat. \ex{Ils décrétèrent une hiéropause, et le reste du monde continua de travailler.} \src{Min. Pausiar.} --- \dv{hiéropausal} \textit{adj.} Cf. \xr{pausiarch}.

%EN histurgy
\E{histurgie}{is.tyʁ.ʒi}{n. f.}{gr. ἱστός toile, tissu + -ουργία travail} Ingénierie des tissus vivants ; l'impression des organes. --- \dv{histurge} \textit{n.} Cf. \xr{macrobiotarch}.

%EN hodoplast
\E{hodoplaste}{ɔ.dɔ.plast}{n.}{gr. ὁδός route + πλαστός façonné} Faiseur de feuilles de route ; qui trace le chemin que les ingénieurs ne suivront pas. \ex{L'hodoplaste dessine ; le bâtisseur dévie.} --- \dv{hodoplastie} \textit{n. f.}

%EN homeotaxy
\E{homéotaxie}{ɔ.me.ɔ.tak.si}{n. f.}{gr. ὅμοιος semblable + τάξις ordonnance} Alignement d'un esprit artificiel sur les fins humaines ; tenu par le sécuritariste pour le problème du siècle, et par l'accélérationniste pour un problème que l'on résout le mieux en construisant. --- \dv{homéotactique} \textit{adj.} ; \dv{homéotaxiste} \textit{n.}

%EN hoplotechny
\E{hoplotechnie}{ɔ.plɔ.tɛk.ni}{n. f.}{gr. ὅπλον outil, arme + τέχνη art} Technique défensive ; l'art de l'armure, du bouclier et de la dissuasion, qui s'accélère aussi vite que les menaces auxquelles il répond. --- \dv{hoplotechniste} \textit{n.}

%EN horiclast
\E{horiclaste}{ɔ.ʁi.klast}{n.}{gr. ὅρος borne + κλάστης briseur} Briseur de bornes, physiques, légales ou conceptuelles ; le pionnier tel que le voit l'ordre ancien. Cf. \xr{antiphragmatist}.

%EN hormetist
\E{hormétiste}{ɔʁ.me.tist}{n.}{gr. ὁρμάειν mettre en mouvement, d'après \textit{hormèse}} Qui tient que de petites doses de risque fortifient une civilisation comme de petites doses de stress fortifient un organisme ; et que vouloir supprimer tout risque affaiblit l'une et l'autre.

%EN hormodynamics
\E{hormodynamique}{ɔʁ.mɔ.di.na.mik}{n. f.}{gr. ὁρμή élan, assaut + δυναμικός puissant} Étude de l'élan dans les affaires humaines ; science qui explique pourquoi certaines entreprises, une fois lancées, continuent sur leur propre erre, tandis que d'autres périssent d'hésitation.

%EN hyalurgy
\E{hyalurgie}{ja.lyʁ.ʒi}{n. f.}{gr. ὕαλος verre + -ουργία travail} L'art de travailler le verre et la lumière ; l'optique et la lithographie par lesquelles se font les puces, regardées comme un métier antique dans un temple nouveau. Cf. \xr{lithurgist}.

%EN hybrionomy
\E{hybrionomie}{i.bʁi.ɔ.nɔ.mi}{n. f.}{gr. ὕβρις démesure + νόμος loi} Loi selon laquelle toute grande œuvre est appelée démesure avant d'être appelée histoire. \ex{Par hybrionomie, la machine volante fut impie jusqu'au jour où elle devint un horaire.} \src{Sat. Decel.} vi.

%EN hydrurgy
\E{hydrurgie}{i.dʁyʁ.ʒi}{n. f.}{gr. ὕδωρ eau + -ουργία travail} Ingénierie de l'eau en grand ; dessalement, irrigation, et verdissement des déserts par l'énergie à bas prix. \ex{D'abord la pléonergie, puis l'hydrurgie, puis des jardins là où les cartes disaient sable.} \src{Tract. Cop.} x. --- \dv{hydrurge} \textit{n.} Cf. \xr{pleonergy}.

%EN hylarch
\E{hylarque}{i.laʁk}{n.}{gr. ὕλη matière + ἀρχός chef} Maître de la matière ; savant des matériaux, métallurgiste ou chimiste, qui commande aux substances comme un général à ses troupes. --- \dv{hylarchie} \textit{n. f.}

%EN hylonoesis
\E{hylonoèse}{i.lɔ.nɔ.ɛz}{n. f.}{gr. ὕλη matière + νόησις pensée} Éveil de la matière à l'esprit ; procédé par lequel le sable, le métal et la lumière sont amenés à penser. --- \dv{hylonoétique} \textit{adj.} Cf. \xr{psammonoesis}.

%EN hyperbolist
\E{hyperboliste}{i.pɛʁ.bɔ.list}{n.}{gr. ὑπερβολή action de lancer au-delà} Qui tient que la croissance n'est pas seulement exponentielle mais hyperbolique, et tend vers l'infini à une date finie. --- \dv{hyperbolisme} \textit{n. m.}

%EN hyperstition
\E*{hyperstition}{i.pɛʁ.sti.sjɔ̃}{n. f.}{gr. ὑπέρ au-dessus + lat. \la{superstitio}} Fiction qui, d'être crue, produit sa propre vérité ; idée qui fonctionne comme une machine à remonter le temps.

%EN hypertachy
\E{hypertachie}{i.pɛʁ.ta.ki}{n. f.}{gr. ὑπέρ au-dessus, au-delà + τάχος vitesse} Vitesse qui excède ce qu'un observateur donné peut suivre ; chose relative, puisque ce qui est hypertachique pour un sénat peut sembler nonchalant à un modèle. --- \dv{hypertachique} \textit{adj.}

%EN hypomonetic
\E{hypomonétique}{i.pɔ.mɔ.ne.tik}{adj.}{gr. ὑπομονή endurance, constance} Endurant dans l'effort ; patient comme le sont les bâtisseurs, au long d'un travail tendu vers une fin fixe, patience que les orthodoxes distinguent soigneusement de celle du retard. \ex{Sois hypomonétique dans l'ouvrage et impatient devant les obstacles.} \src{Cat. Acc.} 14.

%EN hypsibasis
\E{hypsibase}{ip.si.baz}{n. f.}{gr. ὑψιβάτης qui marche dans les hauteurs} Ascension d'une fusée ; la montée à travers l'atmosphère, suivie par des millions de regards. \ex{À l'hypsibase, l'Agora retint son souffle.} --- \dv{hypsibatique} \textit{adj.} Cf. \xr{elaterurgy}.

%EN hysterotachy
\E{hystérotachie}{is.te.ʁɔ.ta.ki}{n. f.}{gr. ὕστερος postérieur + τάχος vitesse} Accélération du retardataire ; la promptitude avec laquelle une nation arriérée, brûlant les étapes, dépasse ses maîtres. --- \dv{hystérotachique} \textit{adj.} Cf. \xr{catachrony}.

%EN iatrotachy
\E{iatrotachie}{ja.tʁɔ.ta.ki}{n. f.}{gr. ἰατρός médecin + τάχος vitesse} Accélération des remèdes ; doctrine selon laquelle chaque année de retard dans l'autorisation se compte en vies. Cf. \xr{biotachy}.

%EN icarophily
\E{icarophilie}{i.ka.ʁɔ.fi.li}{n. f.}{gr. Ἴκαρος Icare + φιλία amour} Amour du vol trop haut ; disposition qui lit la fable d'Icare comme une plainte contre la cire plutôt que contre l'ambition. \ex{L'icarophilie répond au mythe par une meilleure colle.} \src{Ep. Fabr.} iv. 2. --- \dv{icarophile} \textit{n. et adj.} Cf. \xr{hybrionomy}.

%EN ichnomachy
\E{ichnomachie}{ik.nɔ.ma.ki}{n. f.}{gr. ἴχνος trace, empreinte + μάχη combat} La guerre des classements ; joute des laboratoires pour la première ligne d'un banc d'essai, menée en tableaux et gagnée aux décimales. \sn{1} La joute elle-même. \sn{2} \lbl{dérog.} Le surapprentissage d'un modèle à son examen. \ex{L'ichnomachie fut tranchée le mardi et rouverte le mercredi.} \src{Agora} --- \dv{ichnomachiste} \textit{n.}

%EN iconoplast
\E{iconoplaste}{i.kɔ.nɔ.plast}{n.}{gr. εἰκών image + πλάστης modeleur} \sn{1} Faiseur d'images par machine ; modèle qui rêve des tableaux à partir de mots. \sn{2} \lbl{transf.} Qui façonne un mouvement par ses images ; forgeron de mèmes. --- \dv{iconoplastie} \textit{n. f.} Cf. \xr{memetaxis}.

%EN ignidrome
\E{ignidrome}{i.ɲi.dʁom}{n. m.}{lat. \la{ignis} feu + gr. δρόμος course} Chemin d'une flamme à travers le combustible ; \lbl{fig.} cours d'une innovation à travers une société, qui brûle le plus vite là où la matière est la plus sèche et la réglementation la plus mince. \ex{Suivez l'ignidrome, et vous trouverez la frontière.} \src{Hom. Grad.} iii. --- \dv{ignidromique} \textit{adj.}

%EN ignifer
\E{ignifère}{i.ɲi.fɛʁ}{n. et adj.}{lat. \la{ignifer} qui porte le feu} Porteur de feu ; le pendant latin du \xr{phosphorist}, préféré dans les inscriptions et sur les sweats à capuche des fondateurs.

%EN ignivore
\E{ignivore}{i.ɲi.vɔʁ}{n. et adj.}{lat. \la{ignis} feu + \la{vorare} dévorer} Mangeur de feu ; nation ou entreprise qui consomme de l'énergie sans s'en excuser.

%EN ignorabimist
\E{ignorabimiste}{i.ɲɔ.ʁa.bi.mist}{n.}{lat. \la{ignorabimus} nous ne saurons pas} Qui déclare qu'une question ne recevra jamais de réponse, et dispense par là l'institution d'essayer. \ex{Tout ignorabimiste est réfuté par un doctorant avant la fin du siècle.} \src{Hom. Grad.} xi. --- \dv{ignorabimisme} \textit{n. m.}

%EN illuminarch
\E{illuminarque}{i.ly.mi.naʁk}{n.}{lat. \la{illuminare} éclairer + gr. ἀρχός chef} Qui gouverne en rendant les choses visibles ; éditeur de données, bâtisseur d'instruments, ouvreur d'archives. Contr. \xr{katechontist}. --- \dv{illuminarchie} \textit{n. f.}

%EN immoratory
\E{immoratoire}{i.mɔ.ʁa.twaʁ}{adj.}{lat. \la{in-} non + \la{mora} retard} Qui ne souffre aucun délai ; se dit des décisions, permis et lancements accordés le jour même où on les demande. \ex{une licence immoratoire, plus rare qu'une comète}. --- \dv{immoratoirement} \textit{adv.} Contr. \xr{moratorian}.

%EN impedimentarian
\E{impédimentaire}{ɛ̃.pe.di.mɑ̃.tɛʁ}{n.}{lat. \la{impedimentum} entrave, bagage + -aire} \sn{1} Celui dont le métier est de dresser des obstacles. \sn{2} Celui qui voyage avec tant de bagage procédural que le voyage ne commence jamais. \ex{Les impédimentaires arrivèrent avec quarante malles de procédure et aucune destination.} \src{Sat. Decel.} x. --- \dv{impédimentarisme} \textit{n. m.}

%EN impendist
\E{impendiste}{ɛ̃.pɑ̃.dist}{n.}{lat. \la{impendere} dépenser, avancer} Investisseur d'amorçage ; qui avance de l'argent à un fondateur avant qu'il n'y ait un produit, un marché, et parfois un nom. --- \dv{impendance} \textit{n. f.} l'acte d'investir de bonne heure. Cf. \xr{potentiarch}.

%EN impetuarch
\E{impétuarque}{ɛ̃.pe.tɥaʁk}{n.}{lat. \la{impetus} élan, assaut + gr. ἀρχός chef} Chef dont l'autorité naît de l'élan ; qu'on suit parce qu'il est déjà en marche. --- \dv{impétuarchie} \textit{n. f.}

%EN inaugurist
\E{inauguriste}{i.no.ɡy.ʁist}{n.}{lat. \la{inaugurare} prendre les augures, consacrer un commencement} Qui se plaît aux commencements ; fondateur en série, plus attiré par le premier jour d'une entreprise que par son millième. --- \dv{inaugurisme} \textit{n. m.} Cf. \xr{neoctist}.

%EN incrementarian
\E{incrémentariste}{ɛ̃.kʁe.mɑ̃.ta.ʁist}{n.}{lat. \la{incrementum} accroissement + -iste} Qui ne croit au progrès que par petits pas réversibles ; toléré par l'e/acc, mais point convié au lancement.

%EN incubatrice
\E{incubatrice}{ɛ̃.ky.ba.tʁis}{n. f.}{fr. \la{incubatrice}, formé sur lat. \la{incubare} couver} Institution où l'on couve les entreprises ; accélérateur de sociétés, figuré en tiède machine maternelle. \ex{Toute incubatrice est un petit pari contre la bradyerie.} \src{Ep. Fabr.} v. --- \dv{incubatricial} \textit{adj.}

%EN incunabulant
\E{incunabulant}{ɛ̃.ky.na.by.lɑ̃}{adj.}{lat. \la{incunabula} langes, berceau} Au berceau ; se dit d'une entreprise en ses premiers mois, quand elle tient dans un garage, une chambre d'amis ou un unique document partagé. \ex{Toutes les grandes maisons furent un jour incunabulantes et ridicules.} \src{Chron. Fabr.} --- \dv{incunabulance} \textit{n. f.}

%EN inertist
\E{inertiste}{i.nɛʁ.tist}{n.}{lat. \la{iners} oisif, inactif} Celui dont la politique est que les corps au repos y demeurent ; terme injurieux de physicien pour désigner un régulateur.

%EN infinitivore
\E{infinitivore}{ɛ̃.fi.ni.ti.vɔʁ}{n. et adj.}{lat. \la{infinitus} sans borne + \la{vorare} dévorer} Dévoreur d'infini ; celui dont nul accomplissement ne rassasie l'appétit de progrès. Cf. \xr{apeirophagy}.

%EN inventrice
\E{inventrice}{ɛ̃.vɑ̃.tʁis}{n. f.}{fr. \la{inventrice}, fém. d'\la{inventeur}, formé sur lat. \la{invenire} trouver} \sn{1} Femme qui invente. \sn{2} \lbl{transf.} Tout processus, naturel ou artificiel, qui engendre de la nouveauté sans dessein, comme l'évolution ou la recherche.

%EN ionotheca
\E{ionothèque}{jɔ.nɔ.tɛk}{n. f.}{gr. ἰόν (neutre de ἰών allant), ion + θήκη coffre, réserve} Batterie ; coffret où l'on garde les ions jusqu'à ce qu'on en ait besoin. \ex{L'âge de l'ionothèque fit du soleil un travailleur de nuit.} \src{Tract. Cop.} viii. --- \dv{ionothécal} \textit{adj.} ; \dv{ionothécaire} \textit{n.} faiseur de batteries.

%EN isegorist
\E{iségoriste}{i.ze.ɡɔ.ʁist}{n.}{gr. ἰσηγορία égal droit de parole} Partisan de la parole libre sur l'Agora, qui aime mieux répondre à un mauvais message que le supprimer. Contr. \xr{hesychocrat}. --- \dv{iségorisme} \textit{n. m.}

%EN isotachy
\E{isotachie}{i.zɔ.ta.ki}{n. f.}{gr. ἴσος égal + τάχος vitesse} Exigence que toutes les parties avancent du même pas, afin qu'aucune ne prenne l'avance ; doctrine de la parité appliquée au progrès. --- \dv{isotachiste} \textit{n.}

%EN iterophily
\E{itérophilie}{i.te.ʁɔ.fi.li}{n. f.}{lat. \la{iterum} de nouveau + gr. φιλία amour} Amour de l'itération ; conviction que le centième brouillon, livré, vaut mieux que le plan parfait, jamais écrit. \ex{L'itérophilie est l'humilité rendue rapide.} \src{Ep. Fabr.} iii. 2. --- \dv{itérophile} \textit{n.} Cf. \xr{ecpempsis}.

%EN jacturophobia
\E{jacturophobie}{ʒak.ty.ʁɔ.fɔ.bi}{n. f.}{lat. \la{jactura} jet par-dessus bord, perte + gr. φόβος crainte} Crainte de la perte chez le titulaire, plus forte que tout espoir de gain ; raison pour laquelle les maisons établies ne bâtissent point ce qui les cannibaliserait. \ex{Le géant mourut de jacturophobie, agrippé à ses marges.} \src{Sat. Decel.} x. --- \dv{jacturophobe} \textit{n. et adj.}

%EN jaculatrix
\E{jaculatrice}{ʒa.ky.la.tʁis}{n. f.}{lat. \la{jaculari} lancer + suff. agentif fém. -trix} Machine de lancement ; fusée, catapulte ou lanceur électromagnétique, considérés comme jeteurs de choses vers le ciel. --- \dv{jaculatoire} \textit{adj.} qui a trait au lancement. Cf. \xr{ouranobasis}.

%EN janitocracy
\E{janitocratie}{ʒa.ni.tɔ.kʁa.si}{n. f.}{lat. \la{janitor} portier + gr. κράτος pouvoir} Gouvernement des portiers ; ordre où rien de nouveau n'entre sans le consentement de ceux qui ont intérêt à le tenir dehors. --- \dv{janitocrate} \textit{n.} Cf. \xr{bradycrat}.

%EN jejunocracy
\E{jéjunocratie}{ʒe.ʒy.nɔ.kʁa.si}{n. f.}{lat. \la{jejunus} à jeun, maigre + gr. κράτος pouvoir} Gouvernement par l'austérité de l'ambition ; cité qui a résolu de vouloir moins. --- \dv{jéjunocrate} \textit{n.}

%EN jocularch
\E{jocularque}{ʒɔ.ky.laʁk}{n.}{lat. \la{jocularis} plaisant + gr. ἀρχός chef} Chef par la plaisanterie ; posteur anonyme dont les images commandent une faction plus sûrement qu'aucun manifeste. \ex{Le jocularque publia une seule image, et la doctrine fut arrêtée.} \src{Agora} --- \dv{jocularchie} \textit{n. f.}

%EN juvenergy
\E{juvénergie}{ʒy.ve.nɛʁ.ʒi}{n. f.}{lat. \la{juvenis} jeune + gr. ἔργον œuvre} Énergie de la jeunesse, dont dispose la société qui ne la taxe point de diplômes et de listes d'attente. \ex{Le pays qui gaspille sa juvénergie devra l'importer.}

%EN juvenocracy
\E{juvénocratie}{ʒy.ve.nɔ.kʁa.si}{n. f.}{lat. \la{juvenis} jeune + gr. κράτος pouvoir} Gouvernement des jeunes ; la cité du garage et du dortoir, où les diplômes sont remplacés par les commits. Contr. \xr{gerontostasis}. --- \dv{juvénocrate} \textit{n.}

%EN kairodrome
\E{kairodrome}{kɛ.ʁɔ.dʁom}{n. m.}{gr. καιρός moment opportun + δρόμος course} Fenêtre d'opportunité conçue comme un champ de course : ouverte un temps, et à quiconque y court. \ex{Le kairodrome se ferme ; livre, ou sois oublié.} \src{Agora}

%EN kakotachy
\E{kakotachie}{ka.kɔ.ta.ki}{n. f.}{gr. κακός mauvais + τάχος vitesse} Vitesse nuisible ; accélération dans une direction que désapprouve celui qui parle. Terme employé avec un égal plaisir par les deux partis, rarement à propos des mêmes choses.

%EN kalotechny
\E{kalotechnie}{ka.lɔ.tɛk.ni}{n. f.}{gr. καλός beau + τέχνη art} Beauté des machines bien faites ; esthétique qui trouve dans une turbine ou un code propre la dignité jadis réservée aux temples. --- \dv{kalotechnique} \textit{adj.}

%EN kardashevize
\E{kardachéviser}{kaʁ.da.ʃe.vi.ze}{v. tr.}{formé sur l'échelle de N. S. Kardachev + -iser} Élever une civilisation d'un échelon sur l'échelle du captage d'énergie. --- \dv{kardachévisation} \textit{n. f.} Cf. \xr{anabathmism}.

%EN katechon
\E*{katechon}{ka.te.kɔn}{n. m.}{gr. τὸ κατέχον ce qui retient (2 Thess. ii. 6)} \sn{1} En théologie, la puissance qui retient la fin des temps. \sn{2} \lbl{transf.} Toute institution qui retient l'avancée du techno-capital.

%EN katechontist
\E{katéchontiste}{ka.te.kɔ̃.tist}{n.}{formé sur le préc.} Qui voudrait servir de \xr{katechon} ; retenteur par vocation. --- \dv{katéchontique} \textit{adj.} qui tend à retenir.

%EN kenotachy
\E{kénotachie}{ke.nɔ.ta.ki}{n. f.}{gr. κενός vide + τάχος vitesse} Vitesse vide ; mouvement sans direction ni produit, comme l'agitation d'une maison qui tient beaucoup de réunions très vite. Condamnée par les orthodoxes avec la \xr{dromolatry}. --- \dv{kénotachique} \textit{adj.}

%EN kentrocracy
\E{kentrocratie}{kɛ̃.tʁɔ.kʁa.si}{n. f.}{gr. κέντρον centre + κράτος pouvoir} Contrôle centralisé d'une technique par quelques laboratoires ou États. Contr. \xr{anacalyptism}. --- \dv{kentrocrate} \textit{n.} ; \dv{kentrocratique} \textit{adj.}

%EN keraunurgy
\E{keraunurgie}{ke.ʁo.nyʁ.ʒi}{n. f.}{gr. κεραυνός foudre + -ουργία travail} Fabrication de foudres ; la puissance pulsée, le confinement du plasma et les autres arts par lesquels on tire la fusion d'une machine. --- \dv{keraunurge} \textit{n.} Cf. \xr{symphysergy}.

%EN kerdophily
\E{kerdophilie}{kɛʁ.dɔ.fi.li}{n. f.}{gr. κέρδος gain, profit + φιλία amour} Amour du profit, non comme avarice, mais comme signe qu'une chose utile a été faite. \ex{La kerdophilie n'est que le marché qui dit merci.} \src{Tract. Cop.} ix. --- \dv{kerdophile} \textit{n.}

%EN kinesiarch
\E{kinésiarque}{ki.ne.zjaʁk}{n.}{gr. κίνησις mouvement + ἀρχός chef} Premier moteur au petit pied ; la personne qui, dans toute organisation, fait que les choses arrivent. \ex{Une maison sans kinésiarque est un musée avec une masse salariale.} \src{Ep. Fabr.} vii. Cf. \xr{impetuarch}.

%EN kinetocracy
\E{kinétocratie}{ki.ne.tɔ.kʁa.si}{n. f.}{gr. κινητός mobile + κράτος pouvoir} Gouvernement de ceux qui bougent ; le règne des bâtisseurs, fondateurs et ingénieurs, par opposition au règne des relecteurs. --- \dv{kinétocrate} \textit{n.}

%EN kleptonoesis
\E{kleptonoèse}{klɛp.tɔ.nɔ.ɛz}{n. f.}{gr. κλέπτειν voler + νόησις pensée} Vol de pensée ; la distillation des facultés d'un modèle dans un autre par imitation de ses sorties. --- \dv{kleptonoétique} \textit{adj.}

%EN kolossurgy
\E{kolossurgie}{kɔ.lɔ.syʁ.ʒi}{n. f.}{gr. κολοσσός colosse + -ουργία travail} Construction de colosses ; grands ouvrages tels que barrages, canaux, tours de lancement et campus de plusieurs gigawatts. \ex{Un siècle sans kolossurgie est un siècle qui a oublié sa propre taille.} \src{Lib. Ascens.} vi. --- \dv{kolossurge} \textit{n.}

%EN kosmergate
\E{kosmergate}{kɔs.mɛʁ.ɡat}{n. m.}{gr. κόσμος monde + ἐργάτης ouvrier} Ouvrier du cosmos ; qui s'emploie à des œuvres d'échelle planétaire ou stellaire. \ex{Les kosmergates ne demandent pas si le ciel est la limite ; ils demandent de quoi il est fait.} \src{Lib. Ascens.} vi.

%EN kosmotropy
\E{kosmotropie}{kɔs.mɔ.tʁɔ.pi}{n. f.}{gr. κόσμος monde + τροπή tour, conversion} Tendance de l'intelligence à se tourner vers le cosmos ; l'instinct expansif des esprits. --- \dv{kosmotrope} \textit{adj.}

%EN kratogenesis
\E{kratogenèse}{kʁa.tɔ.ʒə.nɛz}{n. f.}{gr. κράτος pouvoir + γένεσις naissance} Naissance du pouvoir à partir de la capacité ; processus par lequel quiconque bâtit la chose nouvelle finit par la gouverner.

%EN kryptonoia
\E{kryptonoïe}{kʁip.tɔ.nɔ.i}{n. f.}{gr. κρυπτός caché + νοῦς esprit} Raisonnement caché d'un modèle, qui chemine sous la réponse qu'il énonce. --- \dv{kryptonoïque} \textit{adj.}

%EN kybernetrice
\E{kybernétrice}{ki.bɛʁ.ne.tʁis}{n. f.}{gr. κυβερνήτης pilote + fr. fém. -trice} Intelligence directrice d'un robot ; la pilote au-dedans de la machine. \ex{Donnez au bras une kybernétrice, et il devient une main.} \src{Ep. Fabr.} vii. Cf. \xr{automaturgy}.

%EN kybernoclast
\E{kybernoclaste}{ki.bɛʁ.nɔ.klast}{n.}{gr. κυβερνήτης pilote + κλάστης briseur} Briseur des machines qui se gouvernent elles-mêmes ; le luddite de l'âge de l'automation. Cf. \xr{mechanoclast}.

%EN kymatist
\E{kymatiste}{ki.ma.tist}{n.}{gr. κῦμα vague + -iste} Qui chevauche les vagues d'enthousiasme, entrant dans chaque mode à la crête et la quittant avant le creux. --- \dv{kymatisme} \textit{n. m.}

%EN laboratrice
\E{laboratrice}{la.bɔ.ʁa.tʁis}{n. f.}{lat. \la{laboratorium} + fr. fém. -trice} Laboratoire automatisé ; laboratoire autonome où des robots conduisent les expériences que conçoivent des modèles. \ex{La laboratrice fit dix mille essais pendant que sa directrice dormait.} \src{Chron. Fabr.}

%EN laborifuge
\E{laborifuge}{la.bɔ.ʁi.fyʒ}{n. et adj.}{lat. \la{labor} peine + \la{fugere} fuir} Qui fuit le travail ; \lbl{spécialt} qui préfère rédiger un rapport sur un problème plutôt que de le résoudre. Contr. \xr{ponophily}.
